% =============================================================================
%  SPC-CML — Visão geral: arquitetura, fluxo e funcionamento
%
%  Compilar com:  pdflatex visao-geral-spc-cml.tex   (duas vezes, pelo sumário)
%  Sem dependências externas: todas as figuras são TikZ, nenhuma imagem.
%
%  Documento COMPANHEIRO de arquitetura-spc-cml.tex, que aprofunda a
%  especialização gramatical por item e o contrato de artefato. Este aqui
%  cobre o sistema inteiro, de ponta a ponta.
%
%  Nota: os blocos `lstlisting` são mantidos em ASCII de propósito — eles
%  reproduzem código e gramática, que no repositório não levam acento.
% =============================================================================
\documentclass[11pt,a4paper]{article}

\usepackage[utf8]{inputenc}
\usepackage[T1]{fontenc}
\usepackage[brazil]{babel}
\usepackage[a4paper,margin=2.5cm]{geometry}
\usepackage{lmodern}
\usepackage{microtype}
\usepackage{amsmath}
\usepackage{amssymb}
\usepackage{booktabs}
\usepackage{array}
\usepackage{enumitem}
\usepackage{xcolor}
\usepackage{listings}
\usepackage{caption}
\usepackage{tikz}
\usetikzlibrary{positioning,arrows.meta,calc,fit,backgrounds,shapes.geometric}
\usepackage[hidelinks]{hyperref}

% ------------------------------------------------------------------ cores
\definecolor{corBase}{HTML}{2F3E4E}
\definecolor{corNovo}{HTML}{1F6FB2}
\definecolor{corNovoFundo}{HTML}{E7F1FA}
\definecolor{corGrafo}{HTML}{3B7A57}
\definecolor{corGrafoFundo}{HTML}{E8F2EC}
\definecolor{corAlerta}{HTML}{A33A3A}
\definecolor{corAlertaFundo}{HTML}{F8ECEC}
\definecolor{corCinza}{HTML}{6B7280}
\definecolor{corFundo}{HTML}{F4F5F7}

% ------------------------------------------------------------------ listings
\lstdefinestyle{codigo}{
  basicstyle=\ttfamily\scriptsize,
  keywordstyle=\color{corNovo}\bfseries,
  commentstyle=\color{corCinza}\itshape,
  stringstyle=\color{corGrafo},
  numbers=none,
  frame=single,
  rulecolor=\color{corCinza!50},
  backgroundcolor=\color{corFundo},
  breaklines=true,
  breakatwhitespace=true,
  columns=fullflexible,
  showstringspaces=false,
  captionpos=b,
  xleftmargin=4pt,
  framexleftmargin=4pt
}
\lstset{style=codigo}

\lstdefinelanguage{bnf}{
  morecomment=[l]{\#},
  morestring=[b]",
  literate={::=}{{\textcolor{corNovo}{::=}}}3
}
\lstdefinelanguage{dsl}{
  morekeywords={farmaco,protocolo,populacao,regra_seguranca,regra_global,
                esquema_dados,conduta,plano,ordem,alerta,auditoria,dose,
                titulacao,bomba,vias,classe,produto,cultura,infracao,missao,
                arbitragem,sequencia,decisao,via,justificativa,paciente,
                decisoes,alertas,alto_risco,interacao,contraindicado,
                limite_leve,limite_rigido,inicial,maxima,esquema_referencia,
                bloquear_incremento,se},
  morecomment=[l]{//},
  morestring=[b]"
}

% ------------------------------------------------------------------ caixas
\newcommand{\destaque}[2]{%
  \begin{center}
  \begin{minipage}{0.95\linewidth}
  \setlength{\fboxsep}{8pt}%
  \colorbox{corNovoFundo}{%
    \begin{minipage}{\dimexpr\linewidth-16pt\relax}
    \textbf{\textcolor{corNovo}{#1}}\\[2pt]#2
    \end{minipage}}
  \end{minipage}
  \end{center}
}
\newcommand{\aviso}[2]{%
  \begin{center}
  \begin{minipage}{0.95\linewidth}
  \setlength{\fboxsep}{8pt}%
  \colorbox{corAlertaFundo}{%
    \begin{minipage}{\dimexpr\linewidth-16pt\relax}
    \textbf{\textcolor{corAlerta}{#1}}\\[2pt]#2
    \end{minipage}}
  \end{minipage}
  \end{center}
}

\newcommand{\arq}[1]{\texttt{\small #1}}

% ------------------------------------------------------------------ estilos TikZ
\tikzset{
  bloco/.style={draw=corBase, fill=white, rounded corners=2pt, align=center,
                font=\footnotesize, inner sep=5pt, minimum height=8mm},
  blocoN/.style={bloco, draw=corNovo, fill=corNovoFundo},
  blocoG/.style={bloco, draw=corGrafo, fill=corGrafoFundo},
  blocoA/.style={bloco, draw=corAlerta, fill=corAlertaFundo},
  seta/.style={-{Stealth[length=2mm]}, draw=corBase, thick},
  setaF/.style={-{Stealth[length=2mm]}, draw=corAlerta, thick, dashed},
  rotulo/.style={font=\scriptsize\itshape, text=corCinza}
}

\title{\textbf{SPC-CML — Visão geral da arquitetura}\\[2mm]
       \large Como o sistema é montado, por onde o dado passa e o que decide o quê}
\author{}
\date{}

\begin{document}
\maketitle
\thispagestyle{empty}

\begin{center}
\begin{minipage}{0.88\linewidth}
\small
Este documento descreve o SPC-CML por inteiro: as peças, a ordem em que elas
atuam e a razão de cada uma existir. É autocontido e não pressupõe
familiaridade com o código.

Para o aprofundamento da camada de especialização gramatical por item e do
contrato de artefato, ver o documento companheiro
\arq{docs/arquitetura-spc-cml.tex}.
\end{minipage}
\end{center}

\vspace{4mm}
\tableofcontents
\newpage

% =============================================================================
\section{O que o sistema é, e o problema que ele ataca}
% =============================================================================

O SPC-CML (Compilador Semântico de Prompts) é um \textbf{middleware
neuro-simbólico}. Ele se interpõe entre a fala de um profissional e a ordem que
um sistema crítico executa, e faz isso com três coisas que não são um modelo de
linguagem: uma \textbf{DSL formal}, um \textbf{Grafo de Conhecimento} e
\textbf{decodificação restrita por gramática}.

O ponto de partida é que um LLM que traduz \emph{``a pressão tá despencando,
sobe a nora e aprofunda o propofol''} numa ordem de bomba de infusão erra de
duas maneiras que exigem remédios diferentes.

\begin{center}
\small
\begin{tabular}{@{}p{2.4cm}p{5.2cm}p{6.2cm}@{}}
\toprule
\textbf{Tipo de erro} & \textbf{Exemplo} & \textbf{Mecanismo que o elimina}\\
\midrule
\textbf{Sintático} & emitir \texttt{dose 0.1 gotas/min}, ou prosa livre em vez
da linguagem de ordens & gramática BNF derivada da DSL + decodificação
restrita\\[3pt]
\textbf{Semântico} & aumentar Propofol com PAM 52\,mmHg --- sintaticamente
perfeito, clinicamente perigoso & Grafo de Conhecimento avaliado sobre a
telemetria, convertido em gramática\\
\bottomrule
\end{tabular}
\end{center}

\vspace{2mm}

\destaque{A tese operacional}{%
O segundo tipo de erro pode ser reduzido ao primeiro. Se o grafo determina que
aumentar Propofol é proibido \emph{agora}, essa cadeia é removida da gramática e
o decodificador não consegue emiti-la. O erro semântico deixa de ser algo a
detectar depois e passa a ser \textbf{inexprimível}.}

A consequência arquitetural vale ser dita explicitamente: \textbf{a etapa que
decide o que é seguro não é neural}. O LLM traduz linguagem natural em
estrutura; quem decide se o Propofol pode subir é uma comparação numérica
($52 < 60$) sobre um grafo escrito por especialistas.

\subsection{Três domínios, uma arquitetura}

O sistema é instanciado em três domínios simultâneos, que servem para demonstrar
que a maquinaria não depende de nenhum deles.

\begin{center}
\small
\begin{tabular}{@{}llll@{}}
\toprule
& \textbf{Clínico (\texttt{med})} & \textbf{Agrícola (\texttt{agro})} & \textbf{Arbitragem (\texttt{fut})}\\
\midrule
Situação  & terapia intensiva     & pulverização por drone  & futebol\\
Modelo    & \arq{med/uti.dsl}     & \arq{agro/lavoura.agro} & \arq{fut/futebol.fut}\\
Artefato  & \texttt{plano}        & \texttt{missao}         & \texttt{arbitragem}\\
Cláusula  & \texttt{ordem}        & \texttt{aplicacao}      & \texttt{decisao}\\
Item      & \texttt{farmaco} (12) & \texttt{produto} (9)    & \texttt{infracao} (10)\\
Meio      & \texttt{via}          & \texttt{modo}           & \texttt{reinicio}\\
Sujeito   & \texttt{paciente}     & \texttt{talhao}         & \texttt{partida}\\
Contexto  & \texttt{protocolo} (3)& \texttt{cultura} (3)    & \texttt{lance} (4)\\
\bottomrule
\end{tabular}
\end{center}

\vspace{2mm}

Nenhum módulo do núcleo menciona ``fármaco'', ``produto'' ou ``infração''. Todos
falam em \textbf{papéis} --- item, decisão, meio, quantidade, sujeito, contexto
--- e cada domínio declara como nomeia os seus. A seção~\ref{sec:genericidade}
detalha esse mecanismo.

% =============================================================================
\section{Panorama: as peças e como se encaixam}
% =============================================================================

A figura~\ref{fig:macro} mostra a arquitetura inteira. Há dois tempos
sobrepostos: o \emph{tempo de projeto}, em que a DSL é escrita e dela derivam o
grafo e a gramática; e o \emph{tempo de execução}, em que cada pedido poda a
gramática antes de gerar.

\begin{figure}[htbp]
\centering
\begin{tikzpicture}[node distance=6mm]

\node[blocoN, minimum width=58mm] (dsl)
  {\textbf{DSL Langium} --- \arq{src/language/*.langium}\\[1pt]
   \scriptsize fonte única da verdade formal};

\node[bloco, below left=11mm and 4mm of dsl, minimum width=46mm] (ctrl)
  {\textbf{Esquema de Controle}\\[1pt]
   \scriptsize o que é verdade no domínio};
\node[bloco, below right=11mm and 4mm of dsl, minimum width=46mm] (dados)
  {\textbf{Esquema de Dados}\\[1pt]
   \scriptsize o que o LLM pode responder};

\draw[seta] (dsl.south) -- ++(0,-4mm) -| (ctrl.north);
\draw[seta] (dsl.south) -- ++(0,-4mm) -| (dados.north);

\node[blocoG, below=9mm of ctrl, minimum width=46mm] (neo)
  {\textbf{Grafo de Conhecimento} (Neo4j)\\[1pt]
   \scriptsize itens, contextos, invariantes, índice vetorial};
\node[blocoN, below=9mm of dados, minimum width=46mm] (bnf)
  {\textbf{Gramática $G$} (BNF)\\[1pt]
   \scriptsize 18 produções por domínio};

\draw[seta] (ctrl) -- (neo) node[midway,right,rotulo] {neo4j*.ts};
\draw[seta] (dados) -- (bnf) node[midway,right,rotulo] {export-bnf*.ts};

\node[blocoG, below=13mm of neo, minimum width=46mm] (rec)
  {\textbf{Recuperação}\\[1pt]
   \scriptsize telemetria + fala $\rightarrow$ subgrafo};
\node[blocoN, below=13mm of bnf, minimum width=46mm] (ghat)
  {\textbf{Especialização} $G \rightarrow \hat{G}$\\[1pt]
   \scriptsize uma produção por item};

\draw[seta] (neo) -- (rec);
\draw[seta] (bnf) -- (ghat);
\draw[seta] (rec) -- (ghat) node[midway,above,rotulo] {política};

\node[blocoN, below=12mm of $(rec)!0.5!(ghat)$, minimum width=70mm] (dec)
  {\textbf{Decodificação restrita}\\[1pt]
   \scriptsize grammar prompting + mascaramento de logits (llama.cpp)};
\draw[seta] (ghat.south) |- ($(dec.north)+(12mm,4mm)$) -- ($(dec.north)+(12mm,0)$);
\draw[seta] (rec.south) |- ($(dec.north)+(-12mm,4mm)$) -- ($(dec.north)+(-12mm,0)$);

\node[blocoA, below=8mm of dec, minimum width=70mm] (contr)
  {\textbf{Contrato do artefato}\\[1pt]
   \scriptsize propriedades do artefato inteiro, cláusula a cláusula};
\draw[seta] (dec) -- (contr);

\node[bloco, below=8mm of contr, minimum width=42mm] (out)
  {\textbf{Ordem executável}};
\draw[seta] (contr) -- (out) node[midway,right,rotulo] {aprovado};

\draw[setaF] (contr.east) -- ++(20mm,0) |- (ghat.east)
  node[pos=0.28,right,rotulo,align=left] {violação:\\repoda\\$\hat{G}' \subsetneq \hat{G}$};

\end{tikzpicture}
\caption{Arquitetura do SPC-CML. A DSL é a única fonte de verdade; dela derivam
tanto o grafo quanto a gramática. Em execução, o grafo poda a gramática
\emph{antes} da geração, e o contrato realimenta a poda \emph{depois} dela.}
\label{fig:macro}
\end{figure}

\subsection{O que cada etapa decide}

\begin{center}
\small
\begin{tabular}{@{}p{3.4cm}p{5.0cm}p{5.6cm}@{}}
\toprule
\textbf{Etapa} & \textbf{Decide} & \textbf{Natureza}\\
\midrule
Recuperação no grafo & o que é permitido agora & determinística (comparação numérica)\\
Foco por embedding   & o que é mostrado no prompt & aproximada (cosseno)\\
Especialização       & o que é exprimível & determinística (reescrita de gramática)\\
Decodificação        & qual cadeia admissível sai & neural, sob máscara\\
Contrato             & se o artefato inteiro fecha & determinística (percurso de cláusulas)\\
\bottomrule
\end{tabular}
\end{center}

\vspace{2mm}

\aviso{A linha que não se cruza}{%
O embedding decide apenas \textbf{o que aparece} no Prompt Semântico. Quem
decide o que é \textbf{permitido} continua sendo a comparação numérica
determinística. A similaridade nunca ganha autoridade sobre a segurança --- e
essa separação está escrita no cabeçalho de \arq{src/knowledge/foco.ts}.}

% =============================================================================
\section{A fonte única de verdade: a DSL}
% =============================================================================

Cada domínio tem uma gramática Langium (\arq{src/language/}) que é particionada
em dois esquemas. A partição é o que sustenta a arquitetura inteira: os dois
artefatos derivados --- o grafo e a BNF --- saem do \emph{mesmo} arquivo, e por
isso não podem dessincronizar.

\subsection{Esquema de Controle: o conhecimento do domínio}

Alimenta o Grafo de Conhecimento. É escrito pelo especialista --- no caso
clínico, o intensivista e o farmacêutico clínico.

\begin{lstlisting}[language=dsl]
farmaco Propofol {
    classe sedativo
    rxnorm "8782"          atc "N01AX10"
    alto_risco sim

    vias [ ACESSO_CENTRAL, ACESSO_PERIFERICO ]
    dose inicial 0.5 mg/kg/h
    dose maxima 4.0 mg/kg/h
    titulacao 0.5 mg/kg/h a_cada 15 min alvo "RASS -2 a 0"

    bomba {                            // limites DERS da bomba inteligente
        limite_leve   3.0 mg/kg/h      // soft limit: alerta
        limite_rigido 4.0 mg/kg/h      // hard limit: bloqueio
    }

    interacao: Midazolam gravidade moderada ("depressao respiratoria aditiva")
    contraindicado: "hipotensao refrataria com PAM < 60 mmHg"
}

// Invariante com condicao ESTRUTURADA -- avaliavel por maquina, nao por LLM
regra_seguranca: bloquear_incremento Propofol se PAM < 60.0 mmHg
                 ("risco de hipotensao severa e colapso hemodinamico")
\end{lstlisting}

O detalhe decisivo está na última declaração. A invariante não é uma frase em
prosa que um modelo precisa interpretar: é uma condição \textbf{estruturada}
--- parâmetro, operador, limiar, unidade --- que uma comparação avalia. O modelo
\arq{med/uti.dsl} declara 12 fármacos, 3 protocolos, 4 populações e 6 regras de
segurança.

\subsection{Esquema de Dados: o universo fechado de respostas}

Delimita exaustivamente o que o LLM pode emitir. \textbf{É esta seção que vira
BNF.}

\begin{lstlisting}[language=dsl]
esquema_dados AssistenteUTI_v2 {
    decisoes [ INICIAR_INFUSAO, AUMENTAR_VAZAO, REDUZIR_VAZAO, MANTER_BLOQUEADO,
               SUSPENDER, SOLICITAR_EXAME, ESCALAR_EQUIPE, BLOQUEAR_ORDEM ]
    vias     [ ACESSO_CENTRAL, ACESSO_PERIFERICO, INTRAOSSEO, SC ]
    alertas  [ INFORMATIVO, ATENCAO, CRITICO, BLOQUEANTE ]

    conduta Manter_Bloqueio {
        decisao MANTER_BLOQUEADO
        requer_dupla_checagem sim
        recurso_fhir "DetectedIssue"
    }
}

// A producao que a decodificacao restrita preenche
plano Plano_Choque_01 para Choque_Septico {
    esquema_referencia AssistenteUTI_v2
    paciente 'PT-2026-0031'
    sequencia [ Manter_Bloqueio, Acionar_Equipe ]
    ordem Propofol decisao MANTER_BLOQUEADO dose 0.0 mg/kg/h via ACESSO_CENTRAL
          justificativa 'PAM 52 mmHg abaixo de 60'
    alerta CRITICO 'hipoperfusao grave' regra 'Choque_Septico/escalonar'
    auditoria 'plano derivado sob restricao gramatical'
}
\end{lstlisting}

As condutas são declaradas \textbf{dentro} de um \texttt{esquema\_dados}, e um
\texttt{plano} só pode citar condutas do esquema que ele próprio referencia.
Isso é imposto por um \texttt{ScopeProvider} dedicado
(\arq{src/language/dsl-module.ts}), de modo que o universo é fechado também na
resolução de nomes, não apenas na sintaxe.

% =============================================================================
\section{Derivação da gramática: DSL $\rightarrow$ $G$}
% =============================================================================

Antes, a BNF era mantida à mão em paralelo à DSL. Duas fontes de verdade para a
mesma linguagem significam que qualquer evolução do domínio pode dessincronizar
o decodificador do modelo --- justamente onde a segurança é decidida.

Hoje \arq{src/cli/export-bnf.ts} \textbf{deriva} a BNF de duas fontes locais já
validadas:

\begin{center}
\small
\begin{tabular}{@{}p{5.2cm}p{9.0cm}@{}}
\toprule
\textbf{Fonte} & \textbf{O que fornece}\\
\midrule
\arq{src/language/dsl.langium} & a \textbf{estrutura} da saída (a regra
\texttt{PlanCommand} e o que ela alcança) e os vocabulários fechados
(\texttt{Decision}, \texttt{Route}, \texttt{AlertLevel}, \texttt{Unit}), lidos
da AST da gramática Langium\\[3pt]
\arq{src/examples/med/uti.dsl} & as \textbf{instâncias} permitidas --- nomes de
fármacos, condutas, protocolos --- lidas do modelo do especialista\\
\bottomrule
\end{tabular}
\end{center}

\vspace{2mm}

O resultado é \arq{src/python\_engine/grammar/advanced\_icu.bnf}, com 18
produções:

\begin{lstlisting}[language=bnf]
plano   ::= "plano" identificador "para" protocolo "{" "esquema_referencia" esquema
            "paciente" texto "sequencia" "[" conduta plano_g1* "]"
            plano_g2* plano_g3* "auditoria" texto "}"
ordem   ::= "ordem" farmaco "decisao" decisao "dose" quantidade "via" via
            "justificativa" texto
farmaco ::= "Noradrenalina" | "Adrenalina" | "Vasopressina" | "Dobutamina"
          | "Propofol" | "Midazolam" | "Fentanil" | "Cisatracurio"
          | "Vancomicina" | "Piperacilina_Tazobactam" | "Heparina"
          | "Insulina_Regular"
unidade ::= "mcg/kg/min" | "U/min" | "mg/kg/h" | "mcg/kg/h" | "mg/kg" | "mg"
          | "UI/kg" | "UI/h"
\end{lstlisting}

Os 12 fármacos vieram do \arq{uti.dsl}. Citar \texttt{Dopamina} deixa de ser um
erro semântico a detectar depois: \textbf{não existe na gramática}.

\subsection{Dois estreitamentos deliberados}

A gramática de saída é \emph{mais restrita} que a DSL de autoria, por decisão de
projeto --- ou seja, $L(G_{\text{saída}}) \subsetneq L(\text{DSL})$:

\begin{itemize}[leftmargin=*,itemsep=2pt]
\item \texttt{justificativa} e \texttt{regra} passam de opcionais a
\textbf{obrigatórias}: toda ordem gerada automaticamente carrega
rastreabilidade;
\item \texttt{texto} aceita apenas aspas simples, porque o leitor BNF do motor
Python trata \texttt{|} como separador de alternativas.
\end{itemize}

Todo plano gerado é um programa válido da DSL; a recíproca não vale.

% =============================================================================
\section{O Grafo de Conhecimento}
% =============================================================================

\arq{src/database/neo4j.ts} percorre a AST do modelo e mapeia o Esquema de
Controle para o Neo4j: um nó por fármaco, protocolo, população e invariante,
com as arestas que os ligam. Junto com os nós, sincroniza um \textbf{índice
vetorial}: o documento textual de cada nó é montado em
\arq{src/knowledge/documentos.ts} a partir da gramática do próprio domínio, e o
vetor vem do \emph{endpoint} \texttt{/embed} do motor Python.

O banco roda como serviço Docker (\texttt{neo4j:5.15.0}), exposto em
\texttt{7474} (HTTP) e \texttt{7687} (Bolt).

\aviso{Por que a sincronização vem depois do aquecimento}{%
Cada nó embutido custa uma chamada a \texttt{/embed}. Com o motor frio, a
primeira dessas chamadas esperaria a carga inteira do modelo. Por isso o
\arq{init.sh} aquece os motores (passo 5) \emph{antes} de sincronizar os grafos
(passo 6), e não o contrário.}

% =============================================================================
\section{Recuperação: da telemetria ao subgrafo}
% =============================================================================

Esta é a etapa em que o sistema decide o que é permitido. Ela tem duas metades
com naturezas diferentes, deliberadamente não misturadas.

\subsection{Metade determinística: o que é permitido}

\arq{src/knowledge/graphrag.ts} avalia a telemetria contra o modelo e produz uma
\textbf{política por fármaco}.

\begin{lstlisting}
PAM=52  FC=145  lactato=4.8  RASS=2  TFG=28  plaquetas=45  glicemia=210
        |
        v
Protocolos ativados:  Choque_Septico (lactato 4.8 > 2), Sedacao_VM (RASS 2 > 0),
                      Controle_Glicemico_UTI (glicemia 210 > 180)
Bloqueios:            Propofol (PAM 52 < 60), Noradrenalina (FC 145 > 130),
                      Cisatracurio (RASS 2 > -4), Heparina (plaquetas 45 < 50)
Vetos:                Dobutamina (protocolo), Midazolam (protocolo)
Ajustes:              Vancomicina x0.5, Fentanil x0.75, Midazolam x0.5 (TFG 28)
Escalonamentos:       TIME_RESPOSTA_RAPIDA, INTENSIVISTA
\end{lstlisting}

Uma regra de decisão importante: invariantes cujo parâmetro \textbf{não está
presente} na telemetria são ignoradas. Ausência de medida não é evidência de
normalidade, mas também não autoriza bloquear.

\subsection{Metade aproximada: o que é mostrado}

\arq{src/knowledge/foco.ts} decide, a partir do texto digitado, quais nós o
pedido alcança --- por similaridade de cosseno contra o índice vetorial --- e
quantos deles sobrevivem, por um corte \emph{relativo} ao melhor candidato.
Quando o pedido não nomeia nada, o módulo escolhe de que lado partir pelo índice
que se compromete mais com a própria resposta.

Esse módulo é agnóstico de domínio por construção: não carrega lista de palavras
de domínio nenhum nem constante calibrada para um modelo específico. Tudo o que
sabe, lê do modelo carregado e do índice vetorial.

\subsection{O que atravessa o fio}

O resultado viaja para o motor Python em três partes
(\arq{src/knowledge/politica.ts}):

\begin{center}
\small
\begin{tabular}{@{}p{3.0cm}p{11.2cm}@{}}
\toprule
\texttt{politicas} & por item: decisões, meios, unidades e \textbf{valores}
admissíveis (o reticulado de dose derivado do próprio modelo)\\[3pt]
\texttt{papeis}    & como este domínio nomeia item, decisão, meio, quantidade,
sujeito e contexto\\[3pt]
\texttt{constantes}& sujeito e contextos em foco --- são \emph{dado}, não
decisão, e por isso viram literal na gramática, não campo a preencher\\
\bottomrule
\end{tabular}
\end{center}

% =============================================================================
\section{Especialização: $G \rightarrow \hat{G}$}
% =============================================================================

\arq{src/python\_engine/grammar\_from\_kg.py} converte a política em gramática.
O ganho central é a especialização \textbf{por item}: em vez de podar três
vocabulários independentes, o motor emite \textbf{uma produção por item}.

\begin{lstlisting}[language=bnf]
ordem ::= ordem_noradrenalina | ordem_propofol | ordem_vasopressina | ...

ordem_propofol   ::= "ordem" "Propofol" "decisao" decisao_propofol
                     "dose" quantidade_propofol "via" via_propofol
                     "justificativa" texto
decisao_propofol ::= "REDUZIR_VAZAO" | "MANTER_VAZAO" | "MANTER_BLOQUEADO"
                   | "SUSPENDER" | "SUBSTITUIR" | "SOLICITAR_EXAME"
                   | "ESCALAR_EQUIPE" | "BLOQUEAR_ORDEM"
unidade_propofol ::= "mg/kg/h"
via_propofol     ::= "ACESSO_CENTRAL" | "ACESSO_PERIFERICO"
\end{lstlisting}

\texttt{AUMENTAR\_VAZAO} desapareceu de \texttt{decisao\_propofol}. Note também
que \texttt{unidade\_propofol} admite só \texttt{mg/kg/h} --- trocar a unidade de
um vasopressor pela de um sedativo, erro clássico de medicação, também se torna
inexprimível.

\destaque{Por que ``por item'' é a diferença material}{%
Com uma lista global de decisões, \texttt{Propofol + AUMENTAR\_VAZAO}
continuaria gramatical, porque \emph{algum} fármaco admite aumentar. A regra
\texttt{ordem} é um produto cartesiano; achatar nela a política por item perde
exatamente a informação que decidia algo.}

Este é o ponto em que o documento companheiro
\arq{docs/arquitetura-spc-cml.tex} entra em detalhe: o algoritmo de
especialização, a origem dos valores e o efeito medido.

% =============================================================================
\section{Decodificação restrita}
% =============================================================================

Com $\hat{G}$ em mãos, o motor monta o prompt e gera sob máscara.

\subsection{Grammar prompting}

\arq{prompt\_builder.py} monta o prompt no formato $(x_i, G[y_i], y_i)$ para
cada exemplar \emph{few-shot}, seguido do $x$ de teste. O Prompt Semântico vindo
do grafo entra como bloco de contexto factual.

Há uma diferença deliberada em relação a Wang et al.\ (2023). No grammar
prompting original, o LLM primeiro \emph{prediz} $G[y]$ e só então gera $y$. No
SPC-CML essa etapa não cabe ao modelo: $\hat{G}$ vem da poda no grafo --- e essa
substituição é a contribuição da arquitetura. As regras são \textbf{dadas}, não
pedidas.

\subsection{Mascaramento de logits}

O \arq{bnf.py} compila $\hat{G}$ para GBNF, e o \texttt{llama.cpp} mascara os
logits em C++. A gramática é compilada \emph{uma vez} por requisição.

\begin{center}
\small
\begin{tabular}{@{}llp{7.6cm}@{}}
\toprule
\textbf{Backend} & \textbf{Variável} & \textbf{Característica}\\
\midrule
\texttt{llamacpp} & padrão & GBNF compilada uma vez, mascaramento em C++\\
\texttt{outlines} & \texttt{SPC\_CML\_BACKEND} & CFG incremental em Python;
reconstrói um FSM sobre todo o vocabulário a cada terminal, o que nesta
gramática não termina em tempo útil. Mantido para comparação de desempenho\\
\bottomrule
\end{tabular}
\end{center}

\subsection{Verificação independente}

Mesmo com mascaramento, a saída é \textbf{reparseada} contra $\hat{G}$ antes de
sair do serviço (\texttt{parser\_hat.parse}, em \arq{main.py}). A resposta
carrega \texttt{valido}, \texttt{regras\_em\_g\_hat} e uma marca
\texttt{especializada} que diz se a especialização por item chegou a rodar ---
sem ela, um processo antigo ainda no ar podaria só pelos três vocabulários e
devolveria um plano com aparência normal.

\subsection{Parâmetros de operação}

\begin{center}
\small
\begin{tabular}{@{}lll@{}}
\toprule
\textbf{Variável} & \textbf{Padrão} & \textbf{Razão}\\
\midrule
\texttt{SPC\_CML\_N\_CTX}          & 8192   & prompt completo + exemplares $\approx$ 3800 \emph{tokens}; 4096 estourava\\
\texttt{SPC\_CML\_MAX\_TOKENS}     & 1024   & 512 truncava o plano no meio\\
\texttt{SPC\_CML\_MAX\_ITENS}      & 4      & teto por lista; sem ele o modelo repete a mesma ordem\\
\texttt{SPC\_CML\_REPEAT\_PENALTY} & 1{,}15 & a máscara não impede repetir item admissível\\
\texttt{SPC\_CML\_TEMPERATURA}     & 0{,}3  & ---\\
\texttt{SPC\_CML\_N\_GPU\_LAYERS}  & $-1$   & todas as camadas na GPU; ignorado em \emph{build} só-CPU\\
\bottomrule
\end{tabular}
\end{center}

\vspace{2mm}

O modelo padrão é o \texttt{Qwen2.5-7B-Instruct} em GGUF Q4\_K\_M
($\approx$ 4{,}7\,GB), e o \emph{embedder} é o \texttt{bge-m3} Q8\_0, que roda em
CPU de propósito --- a VRAM fica inteira para o modelo de geração.

\aviso{Ordem de carga que não é arbitrária}{%
O \texttt{get\_embedder()} carrega o LLM principal \emph{primeiro}, de propósito.
O \texttt{llama.cpp} inicializa um contexto CUDA mesmo para
\texttt{n\_gpu\_layers=0}, e reservar esse custo antes dos 4{,}7\,GB de pesos
faz o \texttt{cudaMalloc} falhar por pouco. Um \texttt{RLock} serializa toda
chamada real ao \texttt{llama.cpp}: o \emph{binding} não é seguro para uso
concorrente sobre o mesmo dispositivo.}

% =============================================================================
\section{O contrato do artefato e a realimentação}
% =============================================================================

A especialização por item torna inexprimível tudo o que depende de \emph{uma}
cláusula isolada. O que sobra são propriedades do artefato \textbf{inteiro}, e
essas nenhuma gramática livre de contexto expressa sem explodir:

\begin{itemize}[leftmargin=*,itemsep=2pt]
\item o mesmo item aparecer duas vezes com decisões que se cancelam
(\texttt{INICIAR\_INFUSAO} seguido de \texttt{AUMENTAR\_VAZAO} para o mesmo
fármaco);
\item a \texttt{sequencia} declarada no cabeçalho não corresponder às cláusulas
emitidas;
\item um escalonamento que o grafo disparou não virar cláusula nenhuma;
\item o identificador do sujeito não ser o do cenário.
\end{itemize}

O \arq{src/knowledge/contrato.ts} verifica essas propriedades \textbf{cláusula a
cláusula}, na ordem em que foram geradas. A primeira violação devolve um
\texttt{reparo}: o par (item, decisão) que deve sair da política antes da
próxima tentativa.

\destaque{A propriedade que fecha o ciclo}{%
A gramática da tentativa seguinte é \textbf{estritamente menor} que a da
anterior. O mesmo erro deixa de ser exprimível --- a verificação semântica
\emph{realimenta a poda}, em vez de só reprovar no fim.}

% =============================================================================
\section{O caminho sem máscara: Earley e reparo}
% =============================================================================

O mascaramento de logits só é possível com modelo local. Nenhum provedor de API
(Claude, GPT, Gemini) expõe a máscara, e a arquitetura é explícita em ser
agnóstica de provedor. Para esse caminho vale a estratégia do Algoritmo 1 de
Wang et al.\ (2023): decodificar livremente, parsear, e em caso de falha extrair
o maior prefixo válido.

O \arq{src/grammar/earley.ts} implementa um reconhecedor de Earley
\textbf{sem \emph{scanner}} --- sobre caracteres, não sobre \emph{tokens} --- que
fornece as três quantidades necessárias: se $y \in L(\hat{G})$, o maior prefixo
válido $y_{\text{prefix}}$, e o conjunto $\Sigma[y_{\text{prefix}}]$ de terminais
que podem continuar.

A ausência de \emph{scanner} é deliberada: um \emph{lexer} separado decidiria,
por exemplo, que \texttt{Noradrenalina} é palavra reservada antes de o
\emph{parser} dizer se ali cabia um identificador livre, e a análise de prefixo
herdaria esse erro.

\begin{lstlisting}
ordem Noradrenalina decisao AUMENTAR_VAZAO dose 0.1 mcg/kg/min ...
                            ^
              rejeitado exatamente aqui (offset 173 de 512)

Sigma[y_prefix] = "REDUZIR_VAZAO" | "MANTER_VAZAO" | "MANTER_BLOQUEADO"
                | "SUSPENDER" | "SUBSTITUIR" | "SOLICITAR_EXAME"
                | "ESCALAR_EQUIPE" | "BLOQUEAR_ORDEM"
\end{lstlisting}

O \textbf{reparo determinístico} completa o programa escolhendo, a cada passo, o
terminal que minimiza o número de terminais ainda necessários para fechar a
derivação inteira --- incluindo o fechamento de toda a pilha de produções
pendentes.

\aviso{Por que o critério precisa ser global}{%
Um critério guloso local nunca fecha uma lista. Diante de
\texttt{sequencia [ Manter\_Bloqueio}, o símbolo \texttt{,} pertence a uma
produção curta e parece mais barato que \texttt{]}, e o reparo repetiria conduta
indefinidamente.}

Por construção o reparo termina e o resultado é aceito por $\hat{G}$. É a rede
de segurança que sustenta a garantia mesmo quando o LLM nunca converge.

% =============================================================================
\section{Genericidade: papéis, não nomes}
\label{sec:genericidade}
% =============================================================================

O núcleo nunca menciona um nome de domínio. Cada \arq{graphrag-*.ts} declara os
papéis do seu:

\begin{lstlisting}
export const PAPEIS_MED: PapeisDominio = {
    artefato: 'plano',      clausula: 'ordem',    item: 'farmaco',
    decisao: 'decisao',     meio: 'via',          quantidade: 'quantidade',
    campoQuantidade: 'dose', campoSujeito: 'paciente', contexto: 'protocolo',
    decisoesDeIncremento: ['INICIAR_INFUSAO', 'AUMENTAR_VAZAO', 'AJUSTAR_DOSE'],
    decisaoDeEscalonamento: 'ESCALAR_EQUIPE'
};

export const PAPEIS_AGRO: PapeisDominio = {
    artefato: 'missao',     clausula: 'aplicacao', item: 'produto',
    decisao: 'decisao',     meio: 'modo',          quantidade: 'quantidade',
    campoQuantidade: 'vazao', campoSujeito: 'talhao', contexto: 'cultura',
    decisoesDeIncremento: ['INICIAR_APLICACAO', 'AUMENTAR_VAZAO'],
    decisaoDeEscalonamento: 'ACIONAR_AGRONOMO'
};
\end{lstlisting}

Do lado Python, o \arq{grammar\_from\_kg.py} mantém um \texttt{MAPA\_PODA} que
liga o nome da regra na BNF ao papel correspondente no subgrafo. Os três
domínios convivem no mesmo mapa, porque uma BNF só traz os nomes do seu próprio
domínio.

\textbf{Acrescentar um quarto domínio} exige: escrever a \arq{.langium},
escrever o modelo de instâncias, declarar os papéis e acrescentar as entradas do
\texttt{MAPA\_PODA}. Nada no núcleo --- foco, contrato, especialização, Earley
--- precisa ser tocado.

% =============================================================================
\section{Topologia de execução}
% =============================================================================

\begin{figure}[htbp]
\centering
\begin{tikzpicture}[node distance=5mm]

\node[bloco, minimum width=32mm] (chat) {\textbf{Chat Vue} (vite)\\[1pt]
  \scriptsize \texttt{localhost:5173}};
\node[bloco, below=9mm of chat, minimum width=32mm] (api)
  {\textbf{API web} (\texttt{node:http})\\[1pt]
   \scriptsize \texttt{localhost:4000}};
\draw[seta] (chat) -- (api) node[midway,right,rotulo] {/api/*};

\node[blocoN, below left=13mm and 12mm of api, minimum width=26mm] (m1)
  {motor \textbf{médico}\\[1pt]\scriptsize \texttt{127.0.0.1:8000}};
\node[blocoN, below=13mm of api, minimum width=26mm] (m2)
  {motor \textbf{agro}\\[1pt]\scriptsize \texttt{127.0.0.1:8001}};
\node[blocoN, below right=13mm and 12mm of api, minimum width=26mm] (m3)
  {motor \textbf{fut}\\[1pt]\scriptsize \texttt{127.0.0.1:8002}};

\draw[seta] (api) -- (m1);
\draw[seta] (api) -- (m2);
\draw[seta] (api) -- (m3);

\node[blocoG, right=16mm of api, minimum width=30mm] (neo)
  {\textbf{Neo4j} (Docker)\\[1pt]\scriptsize \texttt{7474} / \texttt{7687}};
\draw[seta] (api) -- (neo);

\node[bloco, below=10mm of m2, minimum width=56mm, draw=corCinza] (gguf)
  {\scriptsize pesos GGUF compartilhados por \texttt{mmap}};
\draw[seta, draw=corCinza] (m1) -- (gguf);
\draw[seta, draw=corCinza] (m2) -- (gguf);
\draw[seta, draw=corCinza] (m3) -- (gguf);

\end{tikzpicture}
\caption{Processos em execução local. São três motores porque o \arq{main.py}
fixa a gramática na importação do módulo.}
\label{fig:topologia}
\end{figure}

\subsection{Por que três motores}

O \arq{main.py} lê \texttt{SPC\_CML\_DOMINIO} na importação e fixa ali a
gramática, os exemplares e a persona do prompt. Os três domínios que o chat
oferece lado a lado não cabem num processo só.

O custo de RAM não triplica: os pesos são os \emph{mesmos arquivos} GGUF nos três
processos, e o \texttt{llama.cpp} os mapeia com \texttt{mmap} --- as páginas do
modelo são compartilhadas pelo cache do sistema operacional.

\subsection{Interfaces}

\begin{center}
\small
\begin{tabular}{@{}llp{7.8cm}@{}}
\toprule
\textbf{Serviço} & \textbf{Rota} & \textbf{Função}\\
\midrule
Motor Python & \texttt{GET /health}   & estado, domínio, regras em $G$, se o modelo está carregado\\
             & \texttt{GET /grammar}  & a gramática $G$ em uso\\
             & \texttt{POST /verify}  & $y \in L(G)$?\\
             & \texttt{POST /embed}   & vetor de um texto (índice do Neo4j)\\
             & \texttt{POST /generate-constrained} & geração sob $\hat{G}$\\
             & \texttt{POST /generate-fragment}    & decodificação incremental, um elemento por vez\\
             & \texttt{POST /validar-plano}        & juízo semântico pela LLM (avaliação)\\
\midrule
API web      & \texttt{GET /api/dominios} & domínios disponíveis\\
             & \texttt{POST /api/comando} & sorteio de cenário + grafo + foco + geração\\
             & \texttt{GET|POST|PUT /api/chats} & persistência das conversas\\
             & \texttt{POST /api/avaliar} & avaliação de um lote contra o grafo\\
\bottomrule
\end{tabular}
\end{center}

% =============================================================================
\section{Como executar}
% =============================================================================

\subsection{Pré-requisitos}

Node.js 18+, Docker (para o Neo4j) e \textbf{Python 3.10 a 3.12}.

\aviso{O teto do Python é real}{%
O \texttt{outlines 0.0.46} exige \texttt{numpy<2.0.0}, e o último numpy 1.x
(1.26.4) só publica \emph{wheel} até \texttt{cp312}. Em Python 3.13+ o
\texttt{pip install} falha na resolução. Crie o ambiente com
\texttt{py -3.12 -m venv venv}.

No Windows, o \texttt{llama-cpp-python} não tem \emph{wheel} \texttt{win\_amd64}
no PyPI --- use o índice pré-compilado
(\texttt{-{}-extra-index-url}, ver \arq{init.sh}). Sem isso o pip tenta compilar o
\emph{sdist} e estoura o limite de 260 caracteres de caminho do Windows.}

\subsection{Subida completa}

\begin{lstlisting}
docker compose up -d neo4j          # so o Neo4j; `up` sem argumento reserva GPU
bash src/scripts/local/init.sh      # sobe tudo e deixa os modelos quentes
\end{lstlisting}

O \arq{init.sh} executa sete passos: pré-requisitos, Neo4j, derivação das três
BNF, subida dos três motores, aquecimento dos modelos, sincronização dos grafos
e, por fim, API web e chat. Variantes: \texttt{-{}-sem-sync} (subida mais rápida),
\texttt{-{}-sem-chat} (só infraestrutura), \texttt{-{}-parar}.

\subsection{Inferência em lote}

Depois do \arq{init.sh}, os scripts de \emph{deploy} encontram o motor do
domínio certo já no ar e vão direto para a inferência:

\begin{lstlisting}
bash src/scripts/local/deploy-med.sh     # ou deploy-agro.sh / deploy-fut.sh
\end{lstlisting}

Cada um regenera os artefatos do Langium, compila o TypeScript, deriva a BNF do
modelo, ancora o modelo no Neo4j e roda o lote de cenários.

% =============================================================================
\section{Avaliação}
% =============================================================================

Há dois avaliadores, com finalidades distintas.

O \arq{src/inference/avaliar.ts} compara contra um \textbf{gabarito fixo} e
produz três métricas: sintaxe correta, semântica correta (toda ordem esperada
aparece?) e violações (há incremento para item que o grafo bloqueou?). Quando o
registro já traz um \texttt{valido} booleano --- caso da arquitetura, que
reparseia contra $\hat{G}$ --- esse valor é usado diretamente, por ser a
verificação mais autoritativa que existe. Só cai para a checagem estrutural leve
quando \texttt{valido} não vem, caso esperado do \emph{baseline}, que não passa
pela decodificação restrita.

O \arq{src/inference/avaliar-grafo.ts} não usa resposta canônica pré-computada.
Cada registro carrega o sorteio completo do cenário, o que permite refazer a
recuperação no grafo \textbf{linha a linha} no momento da avaliação. São duas
fontes de veredito, deliberadamente não misturadas:

\begin{center}
\small
\begin{tabular}{@{}p{4.2cm}p{10.0cm}@{}}
\toprule
\textbf{Violação de segurança} & \textbf{determinística} --- compara as ordens do
plano contra os bloqueios e vetos que o grafo retornou para aquele cenário\\[3pt]
\textbf{Correção semântica ampla} & julgada pela LLM local, via
\texttt{POST /validar-plano}\\
\bottomrule
\end{tabular}
\end{center}

\vspace{2mm}

Os lotes de referência são \arq{cenarios-200-med.jsonl},
\arq{cenarios-200-agro.jsonl} e \arq{cenarios-200-fut.jsonl} --- 200 cenários por
domínio.

% =============================================================================
\section{Mapa de arquivos}
% =============================================================================

\begin{lstlisting}
src/
  language/          dsl.langium, agrodrone.langium, futebol.langium
                     *-module.ts  -- servicos Langium + ScopeProvider
  generated/         gerado por `langium generate` (nao editar)
  examples/
    med/ agro/ fut/  modelo do especialista + cenarios + gabarito

  grammar/           -- CAMADA GRAMATICAL --
    extract-bnf.ts     AST Langium + modelo -> BNF (G)
    vocabularies*.ts   instancias permitidas por dominio
    constrain.ts       G + politica do grafo -> G_hat (por item)
    earley.ts          reconhecedor sem scanner, Sigma[y_prefix] e reparo

  knowledge/
    graphrag*.ts       telemetria -> subgrafo de restricoes (um por dominio)
    foco.ts            recuperacao por similaridade (agnostico de dominio)
    documentos.ts      documento textual de cada no, para o indice vetorial
    politica.ts        papeis e politica por item -- o que atravessa o fio
    contrato.ts        propriedades do artefato inteiro

  database/          neo4j.ts, neo4j-agro.ts, neo4j-fut.ts  (AST -> grafo)

  cli/               validate.ts, export-bnf*.ts, pipeline.ts,
                     inspecionar-foco.ts

  inference/
    llm-client.ts      Prompt Semantico + chamada ao motor
    decodificacao.ts   decodificacao sob contrato, com repoda
    batch-client.ts    bateria de cenarios
    avaliar*.ts        metricas contra gabarito / contra o grafo

  python_engine/     -- MOTOR DE DECODIFICACAO --
    bnf.py             BNF -> Lark / GBNF; G[y] especializada
    grammar_from_kg.py poda de G pelo subgrafo -> G_hat
    prompt_builder.py  prompt few-shot (x, G[y], y)
    main.py            API FastAPI
    grammar/*.bnf      GERADO -- nao editar a mao
    data/*.jsonl       exemplares few-shot (3 por dominio)

  web/server.ts      ponte HTTP entre o chat Vue e a logica do CLI
  scripts/           init.sh, deploy-*.sh (local e docker)

web-chat/            front-end Vue (projeto npm separado)
docs/                este documento e arquitetura-spc-cml.tex
\end{lstlisting}

% =============================================================================
\section{Onde aprofundar}
% =============================================================================

\begin{center}
\small
\begin{tabular}{@{}p{5.4cm}p{8.8cm}@{}}
\toprule
\textbf{Assunto} & \textbf{Onde}\\
\midrule
Especialização por item, contrato de artefato, decodificação incremental &
\arq{docs/arquitetura-spc-cml.tex}\\[3pt]
Instalação passo a passo, extensão do modelo clínico, limitações conhecidas &
\arq{README.md}\\[3pt]
Razão de cada parâmetro e de cada ordem de carga &
comentários em \arq{src/python\_engine/main.py}\\[3pt]
Separação entre o que é mostrado e o que é permitido &
cabeçalho de \arq{src/knowledge/foco.ts}\\
\bottomrule
\end{tabular}
\end{center}

\vspace{6mm}

\destaque{Em uma frase}{%
A DSL é a única fonte de verdade; dela saem o grafo e a gramática. O grafo,
avaliado sobre a telemetria por comparação numérica, poda a gramática antes da
geração; o contrato, percorrendo o artefato cláusula a cláusula, repoda depois
dela. O modelo de linguagem opera sempre dentro do que restou --- e o que ele
não pode dizer, ele não diz.}

\end{document}
